% Conversão da apresentação da trajetória; origem imutável RELATORIO_ATUAL.odt.
\hypertarget{__RefHeading___Toc4729_3118176633}{}
\hypertarget{_Toc45627141}{}
\refstepcounter{chapter}
% text-h-0001
\chapter*{1 - APRESENTAÇÃO DA TRAJETÓRIA DO ALUNO}
\label{ages-iv:trajectory}
\addcontentsline{toc}{chapter}{1 - APRESENTAÇÃO DA TRAJETÓRIA DO ALUNO}
% text-p-0269
\indent \originaltab{}Minha jornada na Engenharia de Software começou no primeiro semestre de 2023, logo após a conclusão do Ensino Médio. No entanto, o meu envolvimento com o mundo da tecnologia começou muito antes de ingressar na universidade. Sempre fui apaixonado por esse universo, e essa paixão me levou a estudar, por conta própria, diversas tecnologias e linguagens de programação ainda antes de iniciar o curso. Isso facilitou muito a minha adaptação ao ambiente acadêmico, principalmente nos espaços de aprendizado do Prédio 32, onde me senti integrado desde o primeiro momento.\par
% text-p-0270
\indent Logo no meu primeiro semestre, tive a oportunidade de participar da Hackatona 2023, organizada pela PUC-RS, um evento que tive conhecimento por meio da minha coordenadora de curso, Ana Paula Terra Bacelo. O tema da Hackatona era o “Desenvolvimento de uma Inteligência Artificial para um Bem Social", e minha equipe decidiu desenvolver um projeto inovador: uma IA que, utilizando um modelo pré-treinado, fosse capaz de identificar a Língua Brasileira de Sinais (LIBRAS) por meio de uma WebCam e gerar o correspondente output da comunicação. Esse projeto nos rendeu o segundo lugar na competição, que contou com a participação de 14 equipes. Mais do que uma colocação, o projeto abriu portas para que nos tornássemos clientes da Agência Experimental de Engenharia de Software (AGES), onde os alunos de Engenharia de Software, meus colegas de curso, deram continuidade ao desenvolvimento da nossa ideia.\par
% text-p-0271
\indent Hoje, no sétimo semestre do curso, percebo o quanto a minha trajetória ganhou novas proporções. Aquele primeiro contato com a tecnologia, movido pela curiosidade e pela vontade de aprender, deu lugar a uma atuação profissional mais estruturada, na qual assumo responsabilidades como cientista de dados pleno. Ao longo desse caminho, tive a oportunidade de trabalhar com clientes relevantes, incluindo a própria PUCRS, além de participar de projetos junto a órgãos governamentais, nos quais atuei não apenas no desenvolvimento técnico, mas também na coordenação de iniciativas e na tomada de decisões.\par
% text-p-0272
\indent Essa evolução não aconteceu de forma repentina, mas como resultado de um processo contínuo de aprendizado, dedicação e envolvimento com tudo aquilo que me propus a fazer. Hoje, consigo perceber com mais clareza o impacto do meu trabalho, não apenas em termos técnicos, mas também na contribuição que gero para as pessoas e para os contextos em que estou inserido. Entender que o meu papel é importante para que determinadas soluções aconteçam é, sem dúvida, uma das maiores motivações que carrego comigo.\par
% text-p-0273
\indent Ao olhar para toda a minha trajetória, sinto um profundo orgulho do caminho que construí até aqui. Mais do que os resultados alcançados, valorizo a consistência do meu desenvolvimento, as conexões que construí e a forma como venho me posicionando como profissional no mercado. Tenho buscado, diariamente, não apenas evoluir tecnicamente, mas também me tornar alguém que gera valor, que contribui ativamente e que faz com que as coisas aconteçam.\par
% text-p-0274
\indent Tenho a convicção de que ainda há muito a aprender e muitos desafios a enfrentar, mas encaro isso com entusiasmo. Se, no início dessa jornada, havia incertezas sobre o futuro, hoje existe clareza de propósito e confiança no caminho que estou trilhando. E, ao pensar no Eduardo que estava saindo do Ensino Médio, tenho certeza de que ele se sentiria orgulhoso ao ver onde cheguei, e, principalmente, animado com tudo o que ainda está por vir.\par
% text-p-0278 — instrução genérica do modelo omitida.
% text-p-0279 — instrução genérica do modelo omitida.
